\section{Conclusion}
\label{sec:conclusion}

We introduced NestSection, an algorithm for constructing hierarchically
consistent redistricting maps across the congressional, state senate, and state
house chambers.  The central contribution is the \textsc{CompatibleSpines}
algorithm, which identifies the largest geographic structure that all three
chambers can share: the factorization trunk determined by $g = \gcd(C, S, H)$.

The spine compatibility score $\sigma(C, S, H) = (1 - g/\min(C,S,H)) \times 100$
provides a single number summarizing how tightly the three chambers are
structurally linked.  We applied this score to all 50 US states under the 2020
apportionment, finding that 11 states admit strict nesting ($\sigma = 0$) and
39 require best-effort Mode~3 nesting ($\sigma \geq 50$).  The absence of
states in the intermediate range is a structural property of current US
apportionments because it is a structural property of the GCD score itself,
not an artifact of the metric.

The three case studies illustrate the range of outcomes:
\begin{itemize}
  \item \textbf{Oregon} ($g = 6$, $\sigma = 0$): A clean $1:5:10$
        congressional/senate/house ratio admits exact hierarchical nesting.
        The three chambers share a 6-region trunk and diverge only in the
        tail splits within each trunk region.
  \item \textbf{North Carolina} ($g = 2$, $\sigma = 85.7\%$): A single trunk
        bisection is all that can be shared.  Mode~3 nesting with
        a target $\tau_\text{pop} \leq 0.04$ is a measurable design goal, but
        the achieved boundary-zone fraction must be reported by an actual run.
  \item \textbf{Texas} ($g = 1$, $\sigma = 96.8\%$): The primality of the
        31-seat senate makes the three-chamber spine trivial.  Meaningful
        congressional-senate-house nesting would require either approximate
        Mode~3 optimization or restructuring the chamber sizes.
\end{itemize}

The legal argument is constructive: requiring NestSection for compatible states
reduces the degrees of freedom available for partisan optimization before the
map is drawn, not after.  States with $g \geq 2$ (17 states) could be subject
to a best-effort nesting mandate; states with $g \geq 5$ (Oregon and Alabama)
could be subject to a strict nesting mandate with zero boundary tolerance.
Such a statute would be the first anti-gerrymandering rule that operates at
the level of geometric structure rather than partisan outcome, though its legal
force would depend on state law, statutory drafting, and judicial review.

The current implementation provides nesting validation in
\path{bisect-cli::nesting}, suite-level nesting profiles in
\path{bisect-cli::suite}, and audit profile support in \path{rplan-audit}.
Future work will add the full three-chamber constructor, validate it on actual
census geography, and measure the gerrymander-resistance effect empirically.

\paragraph{Acknowledgments.}
The authors thank the NCSL for chamber-size data and the Census Bureau for the
2020 apportionment figures.  The \texttt{bisect} Rust workspace is available
at \url{https://github.com/giodl73-repo/BISECT}.
